% Auto-generated by build_cover.py from journal.toml,
% article.toml and authors.toml. Do not edit by hand --
% edit those files and re-run:  python build_cover.py

% --- Journal ---
\newcommand{\JournalName}{Revista Brasileira de Recursos Hídricos}
\newcommand{\JournalNameEn}{Brazilian Journal of Water Resources}
\newcommand{\JournalAbbrev}{RBRH}
\newcommand{\JournalLogoLeft}{assets/rbrh-logo.jpg}
\newcommand{\JournalLogoRight}{assets/abrhidro-logo.jpg}
\newcommand{\EditorInChief}{Adilson Pinheiro}
\newcommand{\LicenseBadge}{assets/cc-by.pdf}
\newcommand{\LicenseCopyright}{© 2026 The Author(s).}
\newcommand{\LicenseStatement}{This is an Open Access article distributed under the terms of the Creative Commons Attribution License, which permits unrestricted use, distribution, and reproduction in any medium, provided the original work is properly cited.}

% --- Article ---
\newcommand{\DocLanguage}{english}
\newcommand{\IsDraft}{false}
\newcommand{\ArticleType}{Case Study}
\newcommand{\ArticleTitle}{The exceptional hydrological disaster of April--May 2024 in southern Brazil}
\newcommand{\ArticleShortTitle}{RS 2024 flood disaster}
\newcommand{\ArticleCity}{Porto Alegre}
\newcommand{\ArticleDOI}{10.1590/2318-0331.302520240119}
\newcommand{\ArticleVolume}{30}
\newcommand{\ArticleIssue}{2025}
\newcommand{\AssociatedEditor}{Adilson Pinheiro}
\newcommand{\CompetingInterests}{The authors declare no competing interests.}
\newcommand{\DataAvailability}{Data is available upon request.}
\newcommand{\Acknowledgements}{The authors wish to acknowledge CERTEL and ELERA for providing data about hydropower dams. The authors wish to acknowledge HidroSens research group from Universidade Federal de Pelotas for the Patos Lagoon data made available.}
\newcommand{\DatePublished}{17 Mar 2025}
\newcommand{\DateReceived}{05 Nov 2024}
\newcommand{\DateReviewed}{06 Nov 2024}
\newcommand{\DateAccepted}{07 Nov 2024}

% --- Abstracts ---
\newcommand{\AbstractEn}{We present an overview of the unprecedented hydrological event that occurred in Rio Grande do Sul state in late April and early May 2024, causing floods and landslides in a scale that was not seen before in Brazil. The rainfall that caused the disaster is incomparable to any recent event in Brazil at large spatial and temporal scale. The intense and long-lasting precipitation led to record floods from small rivers with drainage areas below 1{,}000~km\textsuperscript{2} to the Patos Lagoon, which drains more than 180{,}000~km\textsuperscript{2}. The rainfall also led to landslides that rank first in Brazil in terms of spatial scale and number of occurrences. Mountain rivers experienced flash floods with extreme water levels, while in the lowlands the large rivers and the Patos Lagoon experienced long duration flooding, with stage remaining high for about 30 days. The flood event is probably ranked first in terms of displaced people and damage to property and infrastructure in Brazil. Almost 200 fatalities occurred due to the disaster and protection structure failures. The April--May hydrological event show a glimpse of what can be expected along the coming decades due to climate change in a region that covers South Brazil, and parts of the neighboring countries.}
\newcommand{\KeywordsEn}{Natural disaster; Extreme hydrological events; Exceptional floods.}
\newcommand{\AbstractPt}{Apresentamos uma visão geral do evento hidrológico sem precedentes que ocorreu no estado do Rio Grande do Sul no final de abril e início de maio de 2024, causando inundações e movimentos de massa em uma escala nunca vista no Brasil. A precipitação que causou o desastre é incomparável a qualquer evento recente no país em termos de escala espacial e temporal. A intensa e prolongada precipitação resultou em enchentes recordes, desde pequenos rios com áreas de drenagem abaixo de 1.000~km\textsuperscript{2} até a Laguna dos Patos, que drena mais de 180.000~km\textsuperscript{2}. As chuvas também levaram a movimentos de massa que ocupam o primeiro lugar no Brasil em termos de extensão espacial e número de ocorrências. Enchentes rápidas com níveis extremos de água ocorreram em regiões montanhosas, enquanto nas áreas mais baixas as cheias tiveram longa duração, com os níveis de água permanecendo elevados por cerca de 30 dias. Este desastre provavelmente ocupa a primeira posição no Brasil em número de pessoas deslocadas e danos a propriedades e infraestruturas. Quase 200 fatalidades ocorreram devido ao desastre e falhas em estruturas de proteção. O evento hidrológico de abril e maio oferece um vislumbre do que pode ser esperado nas próximas décadas devido às mudanças climáticas em uma região que abrange o Sul do Brasil e partes dos países vizinhos.}
\newcommand{\KeywordsPt}{Desastre natural; Eventos hidrológicos extremos; Enchentes excepcionais.}

% --- Authors ---
\newcommand{\AuthorBlock}{Walter Collischonn\,\orcidlink{0000-0002-7630-396X}\textsuperscript{1,*}, Fernando Mainardi Fan\textsuperscript{1}, Iporã Possantti\,\orcidlink{0000-0002-2194-4516}\textsuperscript{1}, Fernando Dornelles\textsuperscript{1}, Rodrigo Paiva\textsuperscript{1}, Matheus Sampaio Medeiros\textsuperscript{1}, Gean Paulo Michel\textsuperscript{1}, Fernando Jorge Corrêa Magalhães Filho\textsuperscript{1}, Sofia Royer Moraes\textsuperscript{2}, Francisco Fernando Noronha Marcuzzo\textsuperscript{3}, Rossano Dalla Lana Michel\textsuperscript{1}, Tamara Leitzke Caldeira Beskow\textsuperscript{4}, Samuel Beskow\textsuperscript{4}, Elisa Helena Leão Fernandes\textsuperscript{5}, Leonardo Laipelt dos Santos\textsuperscript{1}, Anderson Ruhoff\textsuperscript{1}, Masato Kobiyama\textsuperscript{1}, Gilberto Loguercio Collares\textsuperscript{4}, Franco Buffon\textsuperscript{3}, Emanuel Duarte\textsuperscript{3}, Stefany Lima\textsuperscript{1}, Fernando Setembrino Cruz Meirelles\textsuperscript{1}, Daniel Gustavo Allasia Piccilli\textsuperscript{6}}
\newcommand{\AffiliationBlock}{\textsuperscript{1}Universidade Federal do Rio Grande do Sul, Porto Alegre, RS, Brazil; \textsuperscript{2}Universidade do Vale do Taquari, Porto Alegre, RS, Brazil; \textsuperscript{3}Serviço Geológico do Brasil, Porto Alegre, RS, Brazil; \textsuperscript{4}Universidade Federal de Pelotas, Pelotas, RS, Brazil; \textsuperscript{5}Universidade Federal do Rio Grande, Rio Grande, RS, Brazil; \textsuperscript{6}Universidade Federal de Santa Maria, Santa Maria, RS, Brazil}
\newcommand{\CorrespondingBlock}{Walter Collischonn\\ \href{mailto:collischonn@iph.ufrgs.br}{collischonn@iph.ufrgs.br}}

% --- Citations ---
\newcommand{\ArticleCitation}{Collischonn et al. (2025). The exceptional hydrological disaster of April--May 2024 in southern Brazil. Revista Brasileira de Recursos Hídricos, 30, 2025.}
\newcommand{\ArticleShortCitation}{Collischonn et al. (2025). RBRH 30. \url{https://doi.org/10.1590/2318-0331.302520240119}}

% --- Highlights ---
\newcommand{\ArticleHighlights}{%
---\quad An overview of the unprecedented hydrological event that occurred in the state of Rio Grande do Sul.\\
---\quad Record floods ran from small headwater rivers to the Patos Lagoon, which drains more than 180,000~km\textsuperscript{2}.\\
---\quad The April--May hydrological event offers a glimpse of what may be expected in the coming decades due to climate change.%
}
